\chapter{Бариогенез, космическая гладкость и трение Хаббла}

\lectureoverview{Сегодняшнюю лекцию организуют два малых безразмерных числа. Первое — это сохранившийся избыток вещества--антивещества, который в курсе округляется до одной части на \(10^8\). Для его динамического объяснения нужны нарушение барионного числа, асимметрия в паре частица--античастица и отклонение от теплового равновесия: три условия Сахарова. Второе — это первичный контраст плотности порядка \(10^{-5}\) в момент развязки. Гравитация такие возмущения усиливает, а не стирает, поэтому их малость служит мотивировкой инфляции. Затем мы подготовим механику инфляции, выведя уравнение движения однородного скалярного поля в расширяющейся Вселенной и отождествив \(3H\dot\phi\) с трением Хаббла.}

\section{Вопрос об энтропии, поставленный с ног на голову}

Сегодня вечером мы начнём инфляцию, но сначала должны закончить с бариогенезом: образованием избытка вещества над антивеществом. Общая схема надёжнее любой детальной численной модели. Поэтому эту тему можно изложить кратко, хотя её центральное число не выведено из первых принципов.

Исторический путь к этой проблеме начался в 1980 году, когда Боб Вагонер спросил Саваса Димопулоса и Леонарда Сасскинда, почему энтропия Вселенной так велика. В космологической формулировке вопрос состоял в том, почему на каждый сохранившийся протон приходится так много фотонов.

При очень высокой температуре фотоны, кварки, антикварки, электроны и позитроны все были многочисленны и близки к тепловому равновесию. И у кварков, и у антикварков численные плотности были того же общего порядка, что и плотность фотонов. Уже одно это наблюдение подсказывает перевернуть вопрос: почему аннигиляция вообще оставила хоть сколько-нибудь нескомпенсированного вещества?

Приоритет в структурном аргументе принадлежит Андрею Сахарову. Димопулос и Сасскинд независимо пришли к очень близким условиям в 1980 году, а затем, готовя исторический обзор, нашли статью Сахарова 1967 года. Эти условия теперь по праву носят имя Сахарова.

Теперь связь с энтропией можно сформулировать явно. Для излучения чёрного тела энтропия пропорциональна числу фотонов с точностью до фиксированного множителя порядка единицы. Действительно,
\begin{equation}
s_\gamma=\frac{4\rho_\gamma}{3T},
\qquad
n_\gamma\propto T^3,
\end{equation}
поэтому для сопутствующего теплового фотонного газа выполняется \(S_\gamma\propto N_\gamma\). Суть не в численных константах; загадка состоит в огромном отношении числа фотонов к числу протонов.

\begin{classroomqa}
\audiencequestion{Как обилие фотонов связано с энтропией?}
\lecturerresponse{Излучение чёрного тела является тепловым, и его энтропия пропорциональна среднему числу фотонов. В этом контексте вопрос о том, почему космическая энтропия так велика, по существу равносилен вопросу о том, почему в космическом микроволновом фоне на каждый протон приходится так много фотонов. В коэффициент пропорциональности входят численные константы, но для загадки порядка величины они несущественны. \lecturetimestamp{00:04:49--00:05:55}}
\end{classroomqa}

По мере остывания Вселенной частицы и античастицы аннигилировали, тогда как излучение сохранялось и в конечном итоге развязалось. Поэтому удивителен не тот факт, что фотоны выжили; удивительно то, что аннигиляция оставила крошечную нескомпенсированную популяцию вещества.

Полезным безразмерным числом поэтому является
\begin{equation}
\eta_B
\equiv
\frac{N_B-N_{\bar B}}{N_\gamma}
\sim
\frac{N_p-N_{\bar p}}{N_p+N_{\bar p}}
\sim 10^{-8},
\label{eq:c08-baryon-asymmetry}
\end{equation}
где последнее значение — это округлённый порядок, используемый в курсе, а не результат прецизионного определения. Как только возникает число странного вида, ответ ``оно просто с самого начала было таким'' перестаёт быть удовлетворительным завершением рассуждения. Количественный механизм должен объяснить, почему ответ мал и почему у него есть определённый знак.

Называть этот избыток \emph{веществом}, а не \emph{антивеществом}, отчасти условно: веществом называется тот член пары, из которого состоим мы. Физическое утверждение состоит лишь в том, что во всей наблюдаемой области доминирует один знак.

\section{Барионное число и первое условие}

Барионное число нормируется так, что протон или нейтрон несёт \(B=1\). На языке кварков,
\begin{equation}
B=\frac{N_q-N_{\bar q}}{3}.
\label{eq:c08-baryon-number}
\end{equation}
Множитель \(1/3\) существен, потому что каждый обычный барион содержит три валентных кварка. Другие адроны тоже могут нести барионное число, но нестабильные из них в конце концов распадаются; поэтому более общим счётным параметром служит разность числа кварков и антикварков.

\begin{classroomqa}
\classquestion{Следует ли разность числа кварков и антикварков умножать на три или делить на три?}
\lecturerresponse{Её нужно делить на три. Протон содержит три кварка и имеет барионное число единица, так что \(3B=N_q-N_{\bar q}\). Исправление на доске — часть определения, а не изменение физического аргумента. \lecturetimestamp{00:09:23--00:10:17}}
\end{classroomqa}

Барионное число похоже на электрический заряд лишь в том отношении, что тоже может быть сохраняющейся величиной. Электрический заряд также является источником дальнодействующего кулоновского поля; барионное число — нет. Протон создаёт электрическое поле, потому что обладает электрическим зарядом, тогда как нейтрон имеет барионное число, но не имеет суммарного электрического заряда. Следовательно, закон сохранения и роль источника силы — логически разные свойства.

\begin{classroomqa}
\audiencequestion{В чём существенное различие между барионным числом и электрическим зарядом?}
\lecturerresponse{Оба могут служить счётными зарядами, но только электрический заряд, насколько известно, является источником дальнодействующего электромагнитного поля. Барионное число не порождает соответствующей кулоновской силы. \lecturetimestamp{00:11:03--00:12:15}}
\end{classroomqa}

Если бы \(B\) сохранялось в точности, первоначальный избыток оставался бы навсегда замороженным. Рождение пар и аннигиляция могли бы по отдельности изменять \(N_B\) и \(N_{\bar B}\), но только на равные величины, так что их разность не менялась бы. Поэтому динамическая теория, начинающаяся со сбалансированного состояния, не может работать, если барионное число не способно изменяться.

Самый наглядный схематический пример — распад протона. С точки зрения электрического заряда и энергии можно представить
\begin{equation}
p\longrightarrow e^+ +\gamma,
\qquad
\bar p\longrightarrow e^-+\gamma.
\label{eq:c08-decay-pair}
\end{equation}
Эти каналы — модели для учёта нарушения барионного числа, а не предсказания доминирующего экспериментального канала. В них \(B=+1\) или \(-1\) переходит в нуль при сохранении электрического заряда.

\begin{figure}[htbp]
\centering
\begin{coursefitfigure}
\begin{tikzpicture}[x=1cm,y=1cm,>=Latex]
\draw[->,thick] (0,1.7)--(2.0,1.7);
\draw[->,thick] (2.0,1.7)--(4.1,2.55);
\draw[->,thick] (2.0,1.7)--(4.1,0.85);
\node[above] at (0.9,1.7) {\(p,\ B=1\)};
\node[right] at (4.1,2.55) {\(\gamma\)};
\node[right] at (4.1,0.85) {\(e^+\)};
\draw[->,thick] (6.0,1.7)--(8.0,1.7);
\draw[->,thick] (8.0,1.7)--(10.1,2.55);
\draw[->,thick] (8.0,1.7)--(10.1,0.85);
\node[above] at (6.9,1.7) {\(\bar p,\ B=-1\)};
\node[right] at (10.1,2.55) {\(\gamma\)};
\node[right] at (10.1,0.85) {\(e^-\)};
\node at (2.0,0.25) {\(B:1\to0\)};
\node at (8.0,0.25) {\(B:-1\to0\)};
\end{tikzpicture}
\end{coursefitfigure}
\caption{Схематические каналы, связанные преобразованием зарядового сопряжения. Каждый из процессов нарушает барионное число, но равные скорости не создавали бы предпочтительного знака.}
\label{fig:c08-decay-redraw}
\end{figure}

\begin{classroomqa}
\audiencequestion{Мог бы протон вместо этого распасться на позитрон и нейтрино?}
\lecturerresponse{Такое конечное состояние с двумя фермионами не согласуется с фермионной чётностью начального состояния с одним фермионом. Схематическое конечное состояние \(e^++\gamma\) содержит один фермион и достаточно для аргумента о симметрии. Реалистическая теория распада протона может включать другие частицы и каналы, но нынешняя цель — лишь продемонстрировать \(\Delta B\ne0\). \lecturetimestamp{00:14:39--00:15:07}}
\end{classroomqa}

\begin{classroomqa}
\audiencequestion{Соответствовало бы точное сохранение барионного числа некоторой симметрии?}
\lecturerresponse{Да. Сохраняющаяся величина может быть связана с непрерывной фазовой симметрией, хотя одно лишь название этой симметрии само по себе не доказывает, что природа её соблюдает. Операционным критерием остаётся вопрос о том, происходят ли процессы с изменением барионного числа, поэтому рассуждение становится круговым, если симметрию выводят только из закона сохранения. \lecturetimestamp{00:15:08--00:15:46}}
\end{classroomqa}

\begin{classroomqa}
\audiencequestion{Можно ли саму аннигиляцию протона--антипротона рассматривать как распад протона?}
\lecturerresponse{Терминология вторична; закон сохранения — нет. В процессе \(p+\bar p\to\text{photons}\) начальное барионное число равно \(1-1=0\), и конечное барионное число тоже равно нулю. Аннигиляция устраняет две частицы, не изменяя суммарного \(B\), тогда как одночастичный канал в уравнении~\eqref{eq:c08-decay-pair} изменяет \(B\). \lecturetimestamp{00:15:48--00:16:47}}
\end{classroomqa}

Наблюдаемые сегодня протоны чрезвычайно стабильны: в лекции приводятся нижние пределы порядка \(10^{32}\)--\(10^{34}\) лет, что неизмеримо больше возраста Вселенной. Это наблюдение не устанавливает точного сохранения; оно говорит лишь о том, что процессы с изменением барионного числа при нынешних энергиях чрезвычайно медленны.

\section{Дискретные симметрии и необходимость асимметрии}

Следующий вопрос состоит в том, различают ли законы природы вещество и антивещество. Ответ структурируют три дискретных преобразования:
\begin{align}
C&:\ \text{particle}\longleftrightarrow\text{antiparticle},\\
P&:\ \mathbf x\longmapsto-\mathbf x,\\
T&:\ t\longmapsto-t.
\end{align}
Здесь \(P\) меняет пространственную хиральность на противоположную, а \(T\) обращает микроскопический фильм вспять. Макроскопический второй закон не обязан выглядеть инвариантным относительно обращения времени, но вопрос касается лежащей в основе микроскопической динамики.

\begin{figure}[htbp]
\centering
\includegraphics[width=0.82\textwidth]{lecture_08_figure_03.png}
\caption{На доске исправленное определение \(B=(N_q-N_{\bar q})/3\) уже фигурирует при введении \(C\), зарядового сопряжения. Незаконченная запись после слова charge не используется как самостоятельное свидетельство.}
\label{fig:c08-charge-board}
\lectureframe{00:20:23.280}
\end{figure}

Релятивистская локальная квантовая теория поля не требует, чтобы \(C\), \(P\) или \(T\) по отдельности были симметриями. При её стандартных предположениях она требует инвариантности относительно совместного преобразования \(CPT\). Если некоторый процесс возможен, то замена каждой частицы на её античастицу, отражение пространственной конфигурации и обращение движения дают другой допустимый процесс.

Обращение времени в квантовой механике тоньше, чем простая смена знака \(t\). Для бесспиновой волновой функции в потенциале, инвариантном относительно обращения времени,
\begin{equation}
i\hbar\frac{\partial\psi}{\partial t}=H\psi,
\qquad
T:\ \psi(\mathbf x,t)\longmapsto\psi^*(\mathbf x,-t).
\end{equation}
Замена \(t\to-t\) меняет знак производной по времени; комплексное сопряжение также переводит \(i\to-i\), восстанавливая уравнение. Для частиц со спином антиунитарный оператор обращения времени имеет дополнительную структуру, но уже в этом элементарном случае видно комплексное сопряжение.

\begin{classroomqa}
\audiencequestion{Для волновой функции требует ли обращение времени, помимо \(t\to-t\), ещё и комплексного сопряжения?}
\lecturerresponse{Да. Уравнение Шрёдингера показывает, почему: обращение одной лишь производной по времени несогласованно меняет одну сторону, тогда как комплексное сопряжение также меняет знак \(i\). \lecturetimestamp{00:23:58--00:25:04}}
\end{classroomqa}

Теперь можно сформулировать первое условие Сахарова. Если Вселенная начиналась с равными количествами частиц и античастиц, рождение и аннигиляция пар сохраняли бы нулевое суммарное барионное число, если бы не существовало процесса, способного изменять \(B\):
\begin{equation}
\boxed{\text{baryon-number violation}.}
\end{equation}
Без этого симметричное начальное состояние никогда не сможет приобрести ненулевой барионный избыток.

Нынешняя стабильность протона совместима с этим условием. В логике великого объединения нарушение барионного числа может опосредоваться очень тяжёлой частицей массы \(M_X\). При энергиях, намного меньших \(M_X\), эффективная амплитуда схематически подавлена обратными степенями,
\begin{equation}
\begin{aligned}
\mathcal A_{\Delta B\ne0}&\propto \frac{E^2}{M_X^2},\\
\Gamma_p&\sim \frac{m_p^5}{M_X^4}.
\end{aligned}
\label{eq:c08-heavy-suppression}
\end{equation}
Оценка скорости распада имеет обычное для операторов размерности шесть масштабирование, с точностью до констант связи и множителей, зависящих от канала. Очень большое \(M_X\) может сделать холодный протон практически вечным, одновременно допуская реакции с изменением барионного числа в достаточно энергичной ранней плазме.

\begin{classroomqa}
\audiencequestion{Входит ли дополнительная энергия столкновения уже в тяжёлую массу \(M_X\)?}
\lecturerresponse{Нет. \(M_X\) — это масштаб массы покоя переносчика; \(E\) — энергия, доступная в столкновении. Низкоэнергетическое разложение по \(E/M_X\) сильно подавлено при \(E\ll M_X\), а подавление ослабевает, когда тепловой масштаб столкновений приближается к тяжёлому порогу. Буквальный знаменатель, эвристически написанный на доске, следует понимать именно в этом смысле эффективной теории поля. \lecturetimestamp{00:32:32--00:34:02}}
\end{classroomqa}

Предположим, что начальное состояние содержало равные количества частиц и античастиц. Одно лишь нарушение барионного числа всё равно не создавало бы предпочтения. У каждого распада \(p\to e^++\gamma\) был бы зарядово-сопряжённый партнёр \(\bar p\to e^-+\gamma\), а горячая плазма также запускала бы обратные реакции. Равные микроскопические вероятности вызывали бы переходы с изменением барионного числа, но без среднего дрейфа.

\begin{classroomqa}
\audiencequestion{Каково изменение барионного числа на диаграмме с антипротоном?}
\lecturerresponse{У входящего антипротона \(B=-1\), тогда как у электрона и фотона \(B=0\), так что в этом канале барионное число меняется с \(-1\) на нуль. В сопряжённом канале оно меняется с \(+1\) на нуль. Равная частота двух каналов не даёт никакого суммарного предпочтения. \lecturetimestamp{00:37:24--00:37:39}}
\end{classroomqa}

Случайные флуктуации слишком малы, чтобы объяснить наблюдаемый избыток. Если в \(N\) независимых событиях два исхода равновероятны, то
\begin{equation}
\Delta N\sim\sqrt N,
\qquad
\frac{\Delta N}{N}\sim N^{-1/2}.
\end{equation}
Используя грубую оценку из лекции \(N\sim10^{80}\) для числа барионов в наблюдаемой Вселенной, получаем
\begin{equation}
\frac{\Delta N}{N}\sim10^{-40},
\end{equation}
что на тридцать два порядка меньше, чем \(10^{-8}\). Флуктуация типа подбрасывания монеты также не объясняет, почему во всех наблюдаемых областях должен доминировать один и тот же знак.

\begin{classroomqa}
\audiencequestion{Могла ли изначально гораздо большая популяция, скажем \(10^{160}\) пар в сопутствующей области, ставшей нашей наблюдаемой Вселенной, дать флуктуацию порядка \(10^{80}\)?}
\lecturerresponse{Арифметика квадратного корня допускает такое искусственно подобранное число, но оно не объясняет наблюдаемую пространственную когерентность. При неограниченной статистической флуктуации естественно ожидать участки с доминированием вещества и антивещества, а не один знак по всей видимой Вселенной. \lecturetimestamp{00:39:38--00:40:18}}
\end{classroomqa}

Используемый в эпоху лекции наблюдательный аргумент опирается на два вида свидетельств. Обычные ядра, включая гелий, встречаются в космических лучах, тогда как антиядра в обсуждавшихся данных не наблюдались. Одни только антипротоны менее показательны, потому что высокоэнергетические столкновения могут порождать вторичные антипротоны. Кроме того, границы между галактиками из вещества и антивещества создавали бы заметное аннигиляционное излучение. Отсутствие требуемого сигнала от таких границ сильно говорит против равной смеси галактик и антигалактик в наблюдаемой области. Это наблюдательные аргументы в историческом контексте лекции, а не утверждения об областях за пределами нашего горизонта.

\begin{classroomqa}
\audiencequestion{Измеряются ли сами ядра гелия в космических лучах?}
\lecturerresponse{Да. Заряженные ядра гелия измеряются непосредственно. Детальная статистика космических лучей — отдельная специальность, но суть аргумента в контрасте между обычными ядрами и отсутствием сопоставимой популяции антиядер. \lecturetimestamp{00:41:53--00:42:22}}
\end{classroomqa}

Следовательно, нам нужна микроскопическая асимметрия: нарушение \(C\) и, в реалистических обсуждениях бариогенеза, \(CP\). Системы нейтральных \(B\)-мезонов служат лабораторным примером. Нейтральный мезон и его античастица могут распадаться в сопряжённые конечные состояния с неодинаковыми парциальными скоростями,
\begin{equation}
\Gamma(B^0\to f)\ne\Gamma(\bar B^0\to\bar f).
\label{eq:c08-cp-partial}
\end{equation}
Точный канал и численное значение асимметрии на доске не фиксировались; надёжный вывод состоит в экспериментально установленном существовании асимметрии частица--античастица в слабых распадах.

Тем не менее теорема \(CPT\) требует равенства масс и полных ширин распада частицы и её античастицы в вакууме:
\begin{equation}
\Gamma_{\rm tot}(p)=\Gamma_{\rm tot}(\bar p).
\end{equation}
Это не вынуждает совпадать каждую долю ветвления. Различия в одной сопряжённой паре каналов могут компенсироваться различиями в других, например в схематических электронном и мюонном каналах, так что суммы остаются равными.

\begin{classroomqa}
\audiencequestion{Если полные скорости распада протона и антипротона равны, как может возникнуть смещение?}
\lecturerresponse{Равенство полных ширин не означает равенства по каждому каналу. Более того, бариогенез происходит в плазме: в кинетические уравнения входят термические множители заполнения, обратные реакции и каналы рассеяния. Решающее утверждение состоит не в том, что среда просто обходит \(CPT\), а в том, что асимметричные по \(CP\) парциальные процессы могут породить источниковый член, когда детальный баланс утрачен. \lecturetimestamp{00:49:39--00:51:48; 00:53:22--00:57:14}}
\end{classroomqa}

Это и есть второе условие Сахарова:
\begin{equation}
\boxed{C\ \text{and}\ CP\ \text{violation}.}
\end{equation}
Оно придаёт реакциям, изменяющим барионное число, предпочтительное направление. Открытие нарушения \(CP\) незадолго до статьи Сахарова придало этому условию физическую конкретность.

Минимальная Стандартная модель требует аккуратного различения. Она не допускает обычного пертурбативного распада протона, но её электрослабая аномалия допускает непертурбативное нарушение \(B+L\); при высокой температуре сфалеронные переходы могут сделать это нарушение эффективным. Именно таково исправленное содержание позднейшей самопоправки на доске. Вопрос о том, воспроизводят ли одни только источники Стандартной модели космическую асимметрию количественно, — отдельная и гораздо более трудная проблема.

\section{Равновесие и третье условие}

Остаётся ещё одно препятствие. В полном тепловом равновесии каждая реакция уравновешивается обратной. Если барионное число может меняться, соответствующий равновесный химический потенциал не может сохранять произвольное начальное \(B\), и детальный баланс стирает любой возникший избыток. В лекции та же мысль выражается через \(CPT\): у статической равновесной среды нет макроскопической стрелы времени, так что фактор \(T\) не может дать недостающее различие между процессом и его обратным.

Реальная Вселенная расширяется и остывает. Поэтому макроскопическое состояние различает прямое и обратное направление времени. Это существенно только тогда, когда расширение успешно конкурирует с микроскопической релаксацией:
\begin{equation}
\Gamma_{\rm reaction}\lesssim H.
\label{eq:c08-freezeout-condition}
\end{equation}
Если \(\Gamma_{\rm reaction}\gg H\), столкновения быстро восстанавливают равновесие, несмотря на расширение. Если скорость реакции падает ниже хаббловской, плазма уже не может перестроиться, и концентрация вымораживается.

\begin{classroomqa}
\audiencequestion{Почему одного расширения автоматически недостаточно, чтобы вывести систему из равновесия?}
\lecturerresponse{Очень медленно расширяющаяся среда остаётся сколь угодно близкой к равновесию, потому что микроскопические реакции непрерывно её подстраивают. Отклонение становится существенным только тогда, когда время расширения или охлаждения \(H^{-1}\) мало по сравнению с характерным временем реакции \(\Gamma^{-1}\). Тогда по фильму о макроскопическом состоянии уже видно, какое направление является прямым. \lecturetimestamp{01:01:22--01:02:26}}
\end{classroomqa}

Три условия Сахарова теперь можно собрать вместе:
\begin{enumerate}
\item нарушение барионного числа, так что \(B\) может изменяться;
\item \(C\)- и \(CP\)-нарушение, так что это изменение имеет предпочтительный знак;
\item отклонение от теплового равновесия, так что детальный баланс не может скомпенсировать это предпочтение.
\end{enumerate}
Они необходимы для обсуждаемого механизма и задают структурно достаточный путь к избытку. Однако без конкретной модели физики высоких энергий, сети реакций и космологической истории они не определяют его величину.

\begin{classroomqa}
\audiencequestion{Ограничивает ли измеренное значение асимметрии теории высоких энергий?}
\lecturerresponse{В принципе — да. На практике результат зависит от многих неизвестных масс, констант связи, фаз, скоростей реакций и особенностей ранней космологии. Поэтому единственное измеренное число трудно превратить в узкую допустимую область, не выбрав детальной модели. \lecturetimestamp{01:04:20--01:04:52}}
\end{classroomqa}

\begin{classroomqa}
\audiencequestion{Как связаны барионная и электронная асимметрии?}
\lecturerresponse{Электронейтральность связывает их. В схематическом канале удаление протона создаёт позитрон; обращение барионного дисбаланса сдвигает лептонный дисбаланс в компенсирующем направлении. В более общем случае бариогенез и лептогенез связаны сохранением заряда и реакциями Стандартной модели, хотя то, какие именно комбинации сохраняются, зависит от механизма. \lecturetimestamp{01:04:54--01:05:59}}
\end{classroomqa}

\begin{classroomqa}
\audiencequestion{Не слишком ли мало наблюдаемое в Стандартной модели нарушение \(CP\), чтобы объяснить избыток барионов?}
\lecturerresponse{Это стандартный вывод детальных анализов электрослабого бариогенеза, но лекция подчёркивает зависимость расчёта от модели и множество неизвестных на высоких энергиях. Здесь корректно утверждать, что известное низкоэнергетическое нарушение \(CP\) существует, тогда как успешная количественная модель бариогенеза обычно требует дополнительных ингредиентов. \lecturetimestamp{01:06:01--01:07:50}}
\end{classroomqa}

\begin{classroomqa}
\audiencequestion{Не появились ли сегодняшние электроны просто из уничтожения антипротонов?}
\lecturerresponse{Схематические каналы делают такой учёт соблазнительным, но ранняя плазма содержала также независимые, почти симметричные лептонные популяции. Родственные механизмы можно формулировать как бариогенез или лептогенез, а последующие реакции переносят асимметрию между секторами. Историческое слово \emph{бариогенез} не следует принимать за единственную микроскопическую родословную каждого электрона. \lecturetimestamp{01:07:51--01:08:44}}
\end{classroomqa}

\begin{classroomqa}
\audiencequestion{Почему обсуждение обычно ведут в терминах барионов, а не всех лептонов?}
\lecturerresponse{Барионное вещество учитывается по атомным спектрам и другим астрономическим наблюдениям, тогда как космический нейтринный сектор гораздо труднее непосредственно подсчитать. Название отражает как наблюдательную доступность, так и историю этой темы. \lecturetimestamp{01:08:45--01:09:12}}
\end{classroomqa}

\section{Аннигиляция оставляет избыток}

При высокой температуре рождение пар и аннигиляция идут в обоих направлениях,
\begin{equation}
\gamma+\gamma\longleftrightarrow p+\bar p,
\qquad
E_{\gamma\gamma}\gtrsim2m_pc^2\approx2\,\mathrm{GeV}.
\end{equation}
По мере остывания тепловой ванны фотоны, достаточно энергичные для рождения барионной пары, становятся экспоненциально редкими. Аннигиляция остаётся энергетически разрешённой, но обратное рождение уже не может поддерживать равновесие. Почти каждый протон находит антипротон, и именно малый избыток в уравнении~\eqref{eq:c08-baryon-asymmetry} и остаётся. Энергичные фотоны аннигиляции быстро передают свою энергию тепловой ванне.

\begin{classroomqa}
\audiencequestion{Является ли отношение \(10^{-8}\) современной величиной, и когда произошла основная аннигиляция?}
\lecturerresponse{Это используемый в курсе современный порядок величины отношения числа барионов к числу фотонов. Раньше и барионов, и антибарионов было много. Как только обратное рождение пар стало неэффективным, их симметричная компонента аннигилировала и оставила чистый избыток. \lecturetimestamp{01:09:13--01:11:44}}
\end{classroomqa}

\begin{classroomqa}
\audiencequestion{Когда прекратилась аннигиляция протонов--антипротонов?}
\lecturerresponse{Она прекратилась, когда число антипротонов стало пренебрежимо малым. Лекция указывает момент лишь очень приблизительно: при температурах порядка ГэВ, задолго до нуклеосинтеза и, конечно, задолго до рекомбинации. Никакого точного времени из обсуждения не выводится. \lecturetimestamp{01:11:45--01:12:13}}
\end{classroomqa}

Фотоны аннигиляции не исчезли. Они принадлежат истории теплового излучения, итогом которой в конечном счёте является микроволновый фон. Тем не менее сегодня излучение даёт малый вклад в энергетический баланс, потому что
\begin{equation}
\rho_{\rm rad}\propto a^{-4},
\qquad
\rho_m\propto a^{-3}.
\end{equation}
Число фотонов велико, но энергия каждого фотона испытывает красное смещение; материя сохраняет энергию своей массы покоя.

\begin{classroomqa}
\audiencequestion{Если аннигиляция создала так много фотонов, почему сегодня излучение динамически пренебрежимо?}
\lecturerresponse{Эти фотоны всё ещё присутствуют, главным образом как микроволновый фон, но их плотность энергии убывает на одну дополнительную степень масштабного фактора. Расширение и уменьшает их число в единице объёма, и красно смещает каждый фотон. \lecturetimestamp{01:12:14--01:13:18}}
\end{classroomqa}

\begin{classroomqa}
\audiencequestion{Означает ли очень тяжёлый посредник, что Стандартная модель неполна?}
\lecturerresponse{Посредник распада протона в теории великого объединения действительно был бы новой физикой за пределами минимальной Стандартной модели. Позднейшая поправка об аномальном нарушении барионного числа в Стандартной модели не даёт ни такой обычной частицы-посредника распада протона, ни сама по себе полной количественной модели бариогенеза. \lecturetimestamp{01:13:20--01:13:31}}
\end{classroomqa}

\begin{classroomqa}
\audiencequestion{Могло ли сильное неравенство между двумя схематическими каналами распада превратить популяцию, изначально состоящую только из \(p\) и \(\bar p\), в наблюдаемое равенство числа электронов и протонов?}
\lecturerresponse{Можно придумать такой учёт, но горячая равновесная плазма уже содержит большие популяции всех доступных частиц и античастиц. Кроме того, наблюдаемые \(CP\)-асимметрии не дают права произвольно сделать один канал доминирующим. Согласованная модель должна прослеживать эволюцию всей взаимосвязанной сети реакций, а не только двух стрелок на доске. \lecturetimestamp{01:13:33--01:15:10}}
\end{classroomqa}

На этом обсуждение бариогенеза честно заканчивается. Качественные ингредиенты известны; наблюдаемое число не выведено. Более полный ответ требует и более полной теории частиц, и контролируемого описания истории самой ранней Вселенной.

\section{Второе число: первичная неоднородность}

Инфляция отвечает и на другой вопрос, который благодаря космическому микроволновому фону стал количественным: почему Вселенная на больших масштабах столь однородна и изотропна? КМФ имеет необычайно точный спектр чёрного тела с температурой около \(2.7\,\mathrm K\), и эта температура почти одинакова во всех направлениях. То, что когда-то было лишь общим космологическим впечатлением, стало высокоточным фактом.

Гравитация делает эту гладкость ещё более удивительной. Диффузия размазывает каплю чернил, пока её концентрация не станет равномерной. Притягивающая гравитация делает противоположное: область повышенной плотности притягивает окружающее вещество, увеличивая свою плотность и разрежая соседние области. Поэтому малые возмущения, нарастая со временем, становятся зародышами галактик, скоплений и пустот.

\begin{figure}[htbp]
\centering
\begin{coursefitfigure}
\begin{tikzpicture}[x=1cm,y=1cm,>=Latex]
\draw[thick] (-4.2,0) ellipse (1.25 and 0.72);
\draw[very thick] (0,0) ellipse (1.55 and 0.92);
\draw[thick] (4.2,0) ellipse (1.25 and 0.72);
\node at (-4.2,0) {underdense};
\node at (0,0) {overdense};
\node at (4.2,0) {underdense};
\draw[->,thick] (-2.65,0)--(-1.65,0);
\draw[->,thick] (2.65,0)--(1.65,0);
\draw[->,thick] (-2.85,0.75)--(-1.45,0.4);
\draw[->,thick] (2.85,0.75)--(1.45,0.4);
\draw[->,thick] (-2.85,-0.75)--(-1.45,-0.4);
\draw[->,thick] (2.85,-0.75)--(1.45,-0.4);
\end{tikzpicture}
\end{coursefitfigure}
\caption{Гравитационная неустойчивость переносит вещество из менее плотных окрестностей в область повышенной плотности, усиливая начальный контраст, а не размывая его диффузией.}
\label{fig:c08-gravity-growth}
\end{figure}

Если проследить этот рост назад, ранняя Вселенная должна была быть гораздо более гладкой. Точная однородность тоже была бы неприемлема: точная симметрия сохранялась бы, и галактики не могли бы образоваться. Полезная мера — контраст плотности
\begin{equation}
\delta\equiv\frac{\delta\rho}{\rho},
\end{equation}
где \(\delta\rho\) — локальное отклонение от фоновой плотности. Оценки роста структуры, связанные с Jim Peebles и другими ранними космологами, требовали первичной амплитуды порядка
\begin{equation}
\boxed{\frac{\delta\rho}{\rho}\sim10^{-5}}
\end{equation}
к моменту развязки, при этом \(10^{-4}\)--\(10^{-5}\) приводится как исторический диапазон. И здесь снова число делает вопрос более острым, ещё не объясняя своего значения.

В широком смысле инфляционная идея, предложенная Alan Guth в 1980, — это период колоссального, приблизительно экспоненциального расширения. Растягивание небольшой гладкой области на многие порядки величины разбавляет ранее существовавшую кривизну и оттесняет неоднородности на масштабы, превышающие наблюдаемую область. Аналогия со смятым воздушным шаром хорошо передаёт интуицию сглаживания, хотя её не следует понимать как буквальное вложение пространства.

\begin{classroomqa}
\audiencequestion{Как Вселенная может быть изотропной, если вещество поблизости заметно неоднородно в разных направлениях?}
\lecturerresponse{Изотропия — это статистическое утверждение о больших масштабах, а не утверждение о том, что эта комната, Галактика или каждое направление наблюдения выглядят одинаково. Земля локально гориста, но по сравнению со своим радиусом чрезвычайно гладка. Точно так же галактики и скопления — это локальная структура, наложенная на очень гладкий крупномасштабный фон. \lecturetimestamp{01:25:37--01:26:22}}
\end{classroomqa}

\begin{classroomqa}
\audiencequestion{Разве случайный процесс не создавал бы иногда направления с огромным переуплотнением?}
\lecturerresponse{Эти возмущения не являются неограниченными независимыми флуктуациями. Они образуют случайное поле с определённым, почти однородным и изотропным статистическим спектром. Их амплитуды малы, порядка \(10^{-5}\), даже когда их длины волн велики. Спектр и его происхождение — именно то, что ещё предстоит развить, а не предположение о произвольно больших отклонениях, как при подбрасывании монеты. \lecturetimestamp{01:26:26--01:28:23}}
\end{classroomqa}

\begin{classroomqa}
\audiencequestion{Была ли инфляция вызвана высокой температурой, характерной для ранних времен?}
\lecturerresponse{Температура не является главной переменной во время инфляции. Центральное явление — быстрое экспоненциальное расширение, обусловленное тензором энергии-импульса поля. Тепловая ванна снова становится важной, когда инфляция заканчивается и эта энергия превращается в частицы. \lecturetimestamp{01:28:25--01:29:04}}
\end{classroomqa}

\section{Механическая прелюдия: движение с сопротивлением}

Чтобы подготовить уравнения инфляции, рассмотрим камень единичной массы, падающий в вязкой жидкости. Обозначим его вертикальную координату через \(\phi\), предвосхищая скалярное поле, которое вскоре заменит камень. Если его потенциальная энергия равна \(V(\phi)\), то консервативная сила есть
\begin{equation}
F(\phi)=-\frac{dV}{d\phi},
\label{eq:c08-force-potential}
\end{equation}
а без сопротивления уравнение Ньютона имеет вид \(\ddot\phi=F\).

\begin{figure}[htbp]
\centering
\includegraphics[width=0.82\textwidth]{lecture_08_figure_06.png}
\caption{Наглядное соотношение на доске \(F=-dV/d\phi\) для аналогии с падающим камнем. Стрелка вниз задает соглашение о том, что увеличение \(\phi\) направлено вверх.}
\label{fig:c08-force-board}
\lectureframe{01:35:09.513}
\end{figure}

При линейном вязком сопротивлении дополнительная сила равна \(-\gamma\dot\phi\), где \(\gamma>0\). Уравнение принимает вид
\begin{equation}
\ddot\phi+\gamma\dot\phi=F(\phi)
=-V'(\phi).
\label{eq:c08-viscous-motion}
\end{equation}
Сопротивление исчезает в покое и всегда направлено против скорости. Если \(\phi\) возрастает вверх, то у падающего камня \(F<0\) и \(\dot\phi<0\); тогда \(-\gamma\dot\phi>0\), так что сила сопротивления направлена вверх, как и требуется.

\begin{classroomqa}
\classquestion{Почему сила потенциала и сила сопротивления входят с одним и тем же знаком минус? Разве они не должны быть направлены в противоположные стороны?}
\lecturerresponse{Их направления определяются разными величинами. Сила потенциала равна \(-V'(\phi)\). Сопротивление равно \(-\gamma\dot\phi\). При движении вниз и \(F\), и \(\dot\phi\) отрицательны, поэтому \(F\) направлена вниз, а \(-\gamma\dot\phi\) — вверх. Таким образом, алгебра дает противоположные физические направления без изменения ни одной из формул. \lecturetimestamp{01:34:16--01:35:27}}
\end{classroomqa}

По мере того как камень ускоряется, сопротивление растет, пока силы не уравновесятся. При приблизительно постоянной \(F\) установившееся движение удовлетворяет
\begin{equation}
0=F-\gamma\dot\phi_{\rm term},
\qquad
\dot\phi_{\rm term}=\frac{F}{\gamma}.
\end{equation}
Знак правилен: направленная вниз \(F\) дает отрицательную предельную скорость. Поэтому пологий потенциал и большая \(\gamma\) приводят к медленному движению вниз по склону.

\begin{figure}[htbp]
\centering
\begin{tikzpicture}[x=1cm,y=1cm,>=Latex]
\draw[->] (0,0)--(7.0,0) node[right] {\(\phi\)};
\draw[->] (0,0)--(0,4.4) node[above] {\(V(\phi)\)};
\draw[very thick] plot[smooth] coordinates {(0.5,3.7) (1.6,3.55) (2.8,3.2) (4.0,2.55) (5.1,1.65) (6.3,0.6)};
\filldraw (1.45,3.58) circle (0.09);
\draw[->,thick] (1.55,3.48)--(2.25,3.23);
\node[align=left,anchor=west] at (2.45,3.75) {large drag,\\slow downhill motion};
\end{tikzpicture}
\caption{Пологий потенциал при сильном сопротивлении. В полевом аналоге координата камня будет заменена пространственно однородной амплитудой скалярного поля.}
\label{fig:c08-shallow-potential}
\end{figure}

\section{Инфлатон и хаббловское трение}

Инфляция постулирует дополнительное скалярное поле \(\phi\), которое по соглашению называют инфлатоном. Исторически это была смелая гипотеза, введенная для объяснения космической гладкости. Начнем с поля, приблизительно однородного в пространстве, \(\phi=\phi(t)\). Расширение растягивает пространственные неоднородности, тогда как однородная задача выявляет существенную динамику.

Для канонического скалярного поля полная локальная плотность энергии равна
\begin{equation}
\rho_\phi
=
\frac12\dot\phi^{\,2}
+\frac{1}{2a^2}(\boldsymbol\nabla\phi)^2
+V(\phi).
\end{equation}
В однородном приближении \(\boldsymbol\nabla\phi=0\) это сводится к
\begin{equation}
\rho_\phi=\frac12\dot\phi^{\,2}+V(\phi).
\label{eq:c08-scalar-density}
\end{equation}
Первый член — это кинетическая энергия в пространстве конфигураций поля: он измеряет временное изменение, а не поступательное движение в обычном пространстве.

\begin{classroomqa}
\audiencequestion{Означает ли точка над \(\phi\) частную производную по времени?}
\lecturerresponse{Да. Для однородного поля \(\dot\phi=d\phi/dt\). Пространственные градиенты дали бы отдельный положительный член, но при нынешнем предположении они обращаются в нуль. \lecturetimestamp{01:41:08--01:41:34}}
\end{classroomqa}

Проследим за единичной ячейкой в сопутствующих координатах. Ее физический объем равен \(a^3(t)\), поэтому энергия скалярного поля внутри нее равна
\begin{equation}
E_\phi
=a^3(t)
\left[
\frac12\dot\phi^{\,2}+V(\phi)
\right].
\end{equation}
Явный множитель \(a^3(t)\) — новый элемент: физическая ячейка расширяется, хотя ее сопутствующие координаты фиксированы.

\begin{classroomqa}
\audiencequestion{Зависит ли масштабный фактор \(a\) от пространственного положения?}
\lecturerresponse{Не в однородной модели FRW. В данный космический момент времени он одинаков во всем пространстве, но может меняться со временем. Если бы \(\phi\) зависело от положения, пришлось бы восстановить энергию его пространственного градиента. \lecturetimestamp{01:44:12--01:44:49}}
\end{classroomqa}

\begin{classroomqa}
\audiencequestion{Связано ли значение поля с каждой точкой пространства, хотя теперь мы записываем только одно \(\phi(t)\)?}
\lecturerresponse{Да. Поле сопоставляет значение каждой точке пространства-времени. Однородность означает, что в каждый момент времени эти пространственные значения совпадают, сводя динамику к одной общей функции \(\phi(t)\). \lecturetimestamp{01:44:50--01:45:06}}
\end{classroomqa}

\begin{classroomqa}
\audiencequestion{Какое физическое движение описывает кинетический член поля? Поле становится сильнее или слабее?}
\lecturerresponse{Это определяется уравнением движения, начальным значением и скоростью, а также знаком \(V'(\phi)\). Аналогия здесь — частица, направление движения которой нельзя узнать только по одной ее формуле энергии. \lecturetimestamp{01:45:07--01:46:08}}
\end{classroomqa}

Однородный лагранжиан равен разности кинетической и потенциальной энергий,
\begin{equation}
L_\phi=a^3(t)
\left[
\frac12\dot\phi^{\,2}-V(\phi)
\right].
\label{eq:c08-scalar-lagrangian}
\end{equation}
Применяя уравнение Эйлера--Лагранжа, получаем
\begin{align}
\frac{d}{dt}\left(\frac{\partial L_\phi}{\partial\dot\phi}\right)
&=
\frac{\partial L_\phi}{\partial\phi},\\
\frac{d}{dt}\left(a^3\dot\phi\right)
&=-a^3V'(\phi).
\end{align}
Правило произведения дает два члена,
\begin{equation}
a^3\ddot\phi+3a^2\dot a\,\dot\phi
=-a^3V'(\phi).
\end{equation}
Деля на \(a^3\) и определяя \(H=\dot a/a\), получаем
\begin{equation}
\boxed{
\ddot\phi+3H\dot\phi+V'(\phi)=0
}
\label{eq:c08-hubble-friction}
\end{equation}
или, при \(F=-V'(\phi)\),
\begin{equation}
\ddot\phi+3H\dot\phi=F.
\end{equation}
Каждая строка следует из множителя расширяющегося объема в уравнении~\eqref{eq:c08-scalar-lagrangian}.

\begin{classroomqa}
\audiencequestion{Одинаково ли \(H=\dot a/a\) всюду и действительно ли это константа?}
\lecturerresponse{В однородной модели оно одинаково в каждой пространственной точке на срезе космического времени. Оно не обязано быть постоянным во времени; ``постоянная Хаббла'' часто обозначает его современное значение, тогда как \(H(t)\) — это параметр Хаббла на всем протяжении эволюции. \lecturetimestamp{01:50:32--01:50:57}}
\end{classroomqa}

Уравнение~\eqref{eq:c08-hubble-friction} делает аналогию точной. Оно имеет тот же вид, что и вязкое движение, с коэффициентом сопротивления
\begin{equation}
\gamma_{\rm cosmic}=3H.
\end{equation}
Следовательно, расширение замедляет эволюцию однородного скалярного поля. Если \(H\) велико, а потенциал \(V(\phi)\) достаточно полог, инерционный член становится малым по сравнению с членом сопротивления, и поле приближается к соотношению медленного скатывания
\begin{equation}
3H\dot\phi\simeq-V'(\phi).
\end{equation}
Подобно шару, опускающемуся в холодном вязком масле, оно может чрезвычайно долго проходить через пологий потенциал, хотя на него действует сила.

Следующий шаг — подставить плотность энергии скалярного поля в уравнение Фридмана и позволить скорости расширения откликаться на медленно эволюционирующее поле. Эта связанная система порождает инфляцию. Пока же существенная подготовка завершена: поле запасает кинетическую и потенциальную энергию, расширяющийся объем создает хаббловское трение, а медленно скатывающееся поле может сохранять свою потенциальную энергию почти постоянной достаточно долго, чтобы вызывать быстрое расширение.

\section{Что следует запомнить}

Барионная асимметрия и первичный контраст плотности играют параллельные роли. Ни \(10^{-8}\), ни \(10^{-5}\) не объясняются одним лишь называнием, но каждое превращает расплывчатое космологическое впечатление в количественную проблему. Первое ведет к условиям Сахарова; второе подводит к инфляции.

Переход к инфляции столь же точен. Однородное скалярное поле во вселенной FRW подчиняется
\[
\ddot\phi+3H\dot\phi+V'(\phi)=0.
\]
Средний член — не феноменологическая вязкость, введенная вручную. Он следует из роста физического объема \(a^3(t)\). Именно этот результат будет определять ход изложения следующей лекции о вселенной, динамика которой определяется полем, медленно скатывающимся по своему потенциалу.
